% Section 17.6, translated from sv/chapters/17/summary.tex.

\section{Summary}
\label{c17:sec:sammanfattning}

\begin{booktable}{@{}p{8.5em}L@{}}
  \thead{Concept} & \thead{Description} \\
  \midrule
  Phasor & $U = \abs{U}\,\angle\,\delta$ represents $u(t) = \abs{U}\sin(\omega t + \delta)$ \\
  Phasor addition & Convert to rectangular form, add, convert back \\
  Multiplication & $\abs{z_1}\abs{z_2}\,\angle\,(\delta_1 + \delta_2)$ \\
  Division & $(\abs{z_1}/\abs{z_2})\,\angle\,(\delta_1 - \delta_2)$ \\
  Impedance & $Z = U/I$ \\
\end{booktable}

\needspace{6\baselineskip}
\subsection*{What you should be able to do}
After this chapter, you should understand how complex numbers can be used for calculations with
sinusoidal signals.

\testadigsjalv
\begin{enumerate}
  \item A current is given by $i(t) = 5\sin(100\pi t + \pi/3)\,\text{A}$. Give the corresponding
    phasor $I$ in polar form.
  \item Add the phasors $U_1 = 4\,\angle\,0$ and $U_2 = 3\,\angle\,\frac{\pi}{2}$. Give the sum in
    polar form.
\end{enumerate}

\needspace{6\baselineskip}
\subsection*{Next chapter}
\Kapref{18} is a lab in mathematics in C with the standard library \code{math.h}.
